\begin{pastexam}[label={px:ust-108-ee-ds-03}]{臺聯大電機類 108 學年度 第 3 題}
（問答）

(20\%) Illustrate how data changes while using mergesort to sort:\{50 43 51 83 60 33 21\} and write
down merge sort's time complexity.
\end{pastexam}

\begin{pastsoln}
答案：排序結果 $21, 33, 43, 50, 51, 60, 83$；時間複雜度最好、最壞、平均都是 $\Theta(n \lg n)$。

用由上而下的合併排序（資結書虛擬碼~14.5）：$A[p..r]$ 從 $q = \lfloor (p + r)/2 \rfloor$ 切成
$A[p..q]$ 和 $A[q+1..r]$，兩半各自遞迴排好，再合併。 $7$ 個數的位置是 $1..7$，第一刀
$q = \lfloor 8/2 \rfloor = 4$，左半 $4$ 個、右半 $3$ 個。

\textbf{切開。}
\begin{center}
\begin{tabular}{ll}
\toprule
層 & 子陣列 \\
\midrule
0 & $[50\ 43\ 51\ 83\ 60\ 33\ 21]$ \\
1 & $[50\ 43\ 51\ 83]$ \quad $[60\ 33\ 21]$ \\
2 & $[50\ 43]$ \quad $[51\ 83]$ \quad $[60\ 33]$ \quad $[21]$ \\
3 & $[50]$ \quad $[43]$ \quad $[51]$ \quad $[83]$ \quad $[60]$ \quad $[33]$ \\
\bottomrule
\end{tabular}
\end{center}

\textbf{合併（照遞迴實際執行的順序）。} 每一步合併兩段已經排好的子陣列，每次比較兩段最前面的數，
小的先放。 右邊是合併完之後整個陣列的樣子。
\begin{center}
\begin{tabular}{clll}
\toprule
步 & 合併 & 結果 & 整個陣列 \\
\midrule
1 & $[50]$ 和 $[43]$ & $[43\ 50]$ & $43\ 50\ 51\ 83\ 60\ 33\ 21$ \\
2 & $[51]$ 和 $[83]$ & $[51\ 83]$ & $43\ 50\ 51\ 83\ 60\ 33\ 21$ \\
3 & $[43\ 50]$ 和 $[51\ 83]$ & $[43\ 50\ 51\ 83]$ & $43\ 50\ 51\ 83\ 60\ 33\ 21$ \\
4 & $[60]$ 和 $[33]$ & $[33\ 60]$ & $43\ 50\ 51\ 83\ 33\ 60\ 21$ \\
5 & $[33\ 60]$ 和 $[21]$ & $[21\ 33\ 60]$ & $43\ 50\ 51\ 83\ 21\ 33\ 60$ \\
6 & $[43\ 50\ 51\ 83]$ 和 $[21\ 33\ 60]$ & $[21\ 33\ 43\ 50\ 51\ 60\ 83]$ & $21\ 33\ 43\ 50\ 51\ 60\ 83$ \\
\bottomrule
\end{tabular}
\end{center}
第 6 步的比較：$21 < 43$ 放 $21$；$33 < 43$ 放 $33$；$60 > 43$ 放 $43$；$60 > 50$ 放 $50$；
$60 > 51$ 放 $51$；$60 < 83$ 放 $60$；右段用完，把剩下的 $83$ 接上。 共比較 $6$ 次。

如果課本用的是由下而上的合併排序（資結書虛擬碼~14.6，長度 $1, 2, 4, \ldots$ 兩兩合併），過程是
\begin{align*}
&[50][43][51][83][60][33][21] \to [43\ 50][51\ 83][33\ 60][21] \\
&\to [43\ 50\ 51\ 83][21\ 33\ 60] \to [21\ 33\ 43\ 50\ 51\ 60\ 83],
\end{align*}
這組資料剛好和上面每一層的分組一樣，最後結果相同。

\textbf{時間複雜度。} 合併兩段總長 $m$ 的子陣列要 $\Theta(m)$ 時間（每比較一次就放好一個數）。
遞迴式是
\[
T(n) = T(\lceil n/2 \rceil) + T(\lfloor n/2 \rfloor) + \Theta(n),
\]
遞迴樹有 $\lceil \lg n \rceil$ 層，每層合併的總長都是 $n$，所以 $T(n) = \Theta(n \lg n)$。 不管
輸入長什麼樣子，切法都一樣、每層都要把 $n$ 個數搬一次，所以最好、最壞、平均情況都是
$\Theta(n \lg n)$（資結書第~14.4~節）。 合併需要一個長度 $n$ 的暫存陣列，額外空間 $\Theta(n)$。
\end{pastsoln}
